\section{A Guide to the  Diffusion Model Literature}
\label{subsec:guide_to_flow_matching_literature}

There is a whole family of models around diffusion models and flow matching in the literature. When you read these papers, you will likely find a different (but equivalent) way of presenting the material from this class. This makes it sometimes a little confusing to read these papers. For this reason, we want to give a brief overview over various frameworks and their differences and put them also in their historical context. \textbf{This is not necessary to understand the remainder of this document} but rather intended to be a support for you in case you read the literature.

\paragraph{Discrete time vs. continuous time.} The first denoising diffusion model papers \citep{sohl2015deep, song2019generative, ho2020denoising} did not use SDEs but constructed Markov chains in \themebf{discrete time}, i.e. with time steps $t=0,1,2,3,\dots$. To this date, you will find a lot of works in the literature working with this discrete-time formulation. While this construction is appealing due to its simplicity, the disadvantage of the time-discrete approach is that it forces you to choose a time discretization before training. Further, the loss function needs to be approximated via an \themebf{evidence lower bound (ELBO)} - which is, as the name suggests, only a lower bound to the loss we actually want to minimize. Later, \citet{song2020score} showed that these constructions were essentially an approximation of a time-continuous SDEs. Further, the ELBO loss becomes tight (i.e. it is not a lower bound anymore) in the continuous time case (e.g. note that \cref{thm:fm_loss} and \cref{thm:dsm_loss} are equalities and not lower bounds - this would be different in the discrete time case). This made the SDE construction popular because it was considered mathematically "cleaner" and that one could control the simulation error via ODE/SDE samplers post training. It is important to note however that both models employ the same loss and are \textit{not} fundamentally different.

\paragraph{"Forward process" vs probability paths.} The first wave of denoising diffusion models \citep{sohl2015deep, song2019generative, ho2020denoising, song2020score} did not use the term \themebf{probability path} but constructed a noising procedure of a data point $z\in\mathbb{R}^d$ via a so-called \themebf{forward process}. This is an SDE of the form
\begin{align}
\label{eq:forward_process}
    \bar{X}_0=z,\quad \dd \bar{X}_t = \uforw_t(\bar{X}_t)\dd t + \sigforw_t \dd \bar{W}_t
\end{align}
The idea is that after drawing a data point $z\sim \pdata$ one simulates the forward process and thereby corrupts or "noises" the data. The forward process is designed such that for $t\to \infty$ its distribution converges to a Gaussian $\mathcal{N}(0,I_d)$. In other words, for $T\gg 0$ it holds that $\bar{X}_{T}\sim\mathcal{N}(0,I_d)$ approximately. Note that this essentially corresponds to a probability path: the conditional distribution of $\bar{X}_t$ given $\bar{X}_0=z$ is a conditional probability path $\bar{p}_t(\cdot|z)$ and the distribution of $\bar{X}_t$ marginalized over $z\sim \pdata$ corresponds to a marginal probability path $\bar{p}_t$.\footnote{Note however that they use an \themebf{inverted time convention}: $\bar{p}_0(\cdot|z)=\pdata$ here.} However, note that with this construction, we need to know the distribution of $X_t|X_0=z$ in closed form in order to train our models to avoid simulating the SDE. This essentially restrict the vector field $\uforw_t$ to ones such that we know the distribution $\bar{X}_t|\bar{X}_0=z$ in closed form. Therefore, throughout the diffusion model literature, vector fields in forward processes are always of the affine form, i.e. $\uforw_t(x)=a_t x$ for some continuous function $a_t$. For this choice, we can use known formulas of the conditional distribution \citep{sarkka2019applied,song2021sde, karras2022elucidating}:
\begin{align*}
\bar{X}_t|\bar{X}_0=z\sim\mathcal{N}\left(\alpha_t z,\beta_t^2 I\right),\quad\alpha_t=\exp\parr{\int\limits_{0}^{t}a_r\dd r},\quad\beta_t^2=\alpha_t^2\int\limits_{0}^{t}\frac{(\sigforw_r)^2}{\alpha^2_r}dr
\end{align*}
Note that these are simply Gaussian probability paths. Therefore, one can say that \textbf{a forward process is a specific way of constructing a (Gaussian) probability path.} The term probability path was introduced by flow matching \citep{lipman2022flow} to both simplify the construction and make it more general at the same time: First, the "forward process" of diffusion models is never actually simulated (only samples from $\bar{p}_t(\cdot|z)$ are drawn during training). Second, a forward process only converges for $t\to\infty$ (i.e. we will never arrive at $\pinit$ in finite time). Therefore, we choose to use probability paths in this document.

\paragraph{Time-Reversals vs Solving the Fokker-Planck equation.}

The original description of diffusion models did not construct the training target $\uref_t$ or $\nabla\log p_t$ via the Fokker-Planck equation (or Continuity equation) but via a \themebf{time-reversal} of the forward process \citep{anderson1982reverse}. A time-reversal $(X_t)_{0\leq t\leq T}$ is an SDE with the same distribution over trajectories inverted in time, i.e.
\begin{align}
\mathbb{P}[\bar{X}_{t_1}\in A_1,\dots,\bar{X}_{t_n}\in A_n]=\mathbb{P}[X_{T-t_1}\in A_1,\dots,X_{T-t_n}\in A_n]\\
    \text{ for all }0\leq t_1,\dots, t_n\leq T, \text{ and } A_1,\dots,A_n\subset S
\end{align}
As shown in \citet{anderson1982reverse}, one can obtain a time-reversal satisfying the above condition by the SDE:
\begin{align*}
    \dd X_t =& \left[-u_t(X_t)+\sigma_t^2\nabla\log p_t(X_t)\right]\dd t+ \sigma_{t}\dd W_t,\quad u_t(x)=\uforw_{T-t}(x),\sigma_t=\bar{\sigma}_{T-t}
\end{align*}
As $u_t(X_t)=a_tX_t$, the above corresponds to a specific instance of training target we derived in \cref{prop:conversion_formula_gaussian_prob_path} (this is not immediately trivial as different time conventions are used. See e.g. \citep{lipman2024flow} for a derivation). However, for the purposes of generative modeling, we often only use the final point $X_1$ of the Markov process (e.g., as a generated image) and discard earlier time points. Therefore, whether a Markov process is a ``true'' time-reversal or follows along a probability path does not matter for many applications. Therefore, using a time-reversal is not necessary and often leads to suboptimal results, e.g. the probability flow ODE is often better \citep{karras2022elucidating, ma2024sit}. All ways of sampling from a diffusion models that are different from the time-reversal rely again on using the Fokker-Planck equation. We hope that this illustrates why nowadays many people construct the training targets directly via the Fokker-Planck equation - as pioneered by \citep{lipman2022flow,liu2022flow,albergo2023stochastic} and done in this class.

\paragraph{Flow Matching \citep{lipman2022flow} and Stochastic Interpolants \citep{albergo2023stochastic}.}  The framework that we present is most closely related to the frameworks of flow matching and \themebf{stochastic interpolants (SIs)}. As we learnt, flow matching restricts itself to flows. In fact, one of the key innovations of flow matching was to show that one does not need a construction via a forward process and SDEs but flow models alone can be trained in a scalable manner. Due to this restriction, you should keep in mind that sampling from a flow matching model will be  deterministic (only the initial $X_0\sim \pinit$ will be random). Stochastic interpolants included both the pure flow and the SDE extension via "Langevin dynamics" that we use here (see \cref{thm:langevin_trick}). Stochastic interpolants get their name from a \themebf{interpolant function} $I(t,x,z)$ intended to interpolate between two distributions. In the terminology we use here, this corresponds to a different yet (mainly) equivalent way of constructing a conditional and marginal probability path. The advantage of flow matching and stochastic interpolants over diffusion models is both their simplicity and their generality: their training framework is very simple but at the same time they allow you to go from an arbitrary distribution $\pinit$ to an arbitrary distribution $\pdata$ - while denoising diffusion models only work for Gaussian initial distributions and  Gaussian probability path. This opens up new possibilities for generative modeling that we will touch upon briefly later in this class.

\begin{summarybox}[Alternative Diffusion Formulations]
Alternative formulations for diffusion models that are popular in the literature often involve some combination of the following elements:
\begin{enumerate}
    \item \textbf{Discrete-time: }Approximations of SDEs via discrete-time Markov chains are often used.
    \item \textbf{Inverted time convention: }It is popular to use an inverted time convention where $t=0$ corresponds to $\pdata$ (as opposed to here where $t=0$ corresponds to $\pinit$).
    \item \textbf{Forward process: }Forward processes (or noising processes) are ways of constructing (Gaussian) probability paths.
    \item \textbf{Training target via time-reversal:} A training target can also be  constructed via the time-reversal of SDEs. This is a specific instance of the construction presented here (with an inverted time convention).
\end{enumerate}

\end{summarybox}